\begin{thebibliography}{99}
\bibitem{bridge:Lewin}
L.~Lewin, \emph{Polylogarithms and Associated Functions}, North-Holland, 1981; and (ed.)
\emph{Structural Properties of Polylogarithms}, AMS, 1991. (The trigamma Clausen-DFT law;
cited via MathWorld, \emph{Trigamma Function}, eq.~(14).)

\bibitem{bridge:WolframFunctions}
Wolfram Research, \emph{The Wolfram Functions Site}, \texttt{functions.wolfram.com},
PolyLog identity 10.08.03.0007.01. (The forward cyclotomic formula.)

\bibitem{bridge:Kummer}
E.~E.~Kummer, \emph{Beitrag zur Theorie der Function
$\Gamma(x)$}, J.~reine angew.\ Math.\ \textbf{35} (1847), 1--4. (The Fourier series of
$\log\Gamma$.)

\bibitem{bridge:Malmsten}
C.~J.~Malmsten, \emph{De integralibus quibusdam definitis, seriebusque infinitis},
J.~reine angew.\ Math.\ \textbf{38} (1849), 1--39 (dated 1846). (The log-log integrals and
the reflection identity for $\gamma_1$.)

\bibitem{bridge:Vardi}
I.~Vardi, \emph{Integrals, an introduction to analytic number theory}, Amer.\ Math.\
Monthly \textbf{95} (1988), 308--315.

\bibitem{bridge:BlagouchineJNT}
I.~V.~Blagouchine, \emph{A theorem for the closed-form evaluation of the first generalized
Stieltjes constant at rational arguments and some related summations}, J.~Number Theory
\textbf{148} (2015), 537--592. (Vendored in the repository at
\texttt{src/PolyLog/docs/pdfs/1401.3724v3/}.)

\bibitem{bridge:EspinosaMoll}
O.~Espinosa and V.~Moll, \emph{On some integrals involving the Hurwitz zeta function,
Part~1}, Ramanujan J.\ \textbf{6} (2002), 159--188. (The second log-gamma moment.)

\bibitem{bridge:WhittakerWatson}
E.~T.~Whittaker and G.~N.~Watson, \emph{A Course of Modern Analysis}, 4th ed., Cambridge
University Press, 1927. (Binet's formulas; the Hurwitz-series ladder.)

\bibitem{bridge:Deninger}
C.~Deninger, \emph{On the analogue of the formula of Chowla and Selberg for real quadratic
fields}, J.~reine angew.\ Math.\ \textbf{351} (1984), 171--191. (The $\zeta''(0,a)$
layer.)

\bibitem{gamma:KoblitzOgus}
N.~Koblitz and A.~Ogus, \emph{Algebraicity of some products of values of the $\Gamma$
function}, appendix (pp.~343--346) to P.~Deligne, \emph{Valeurs de fonctions $L$ et
p\'eriodes d'int\'egrales}, in \emph{Automorphic Forms, Representations and
$L$-Functions} (Corvallis, 1977), Proc.\ Sympos.\ Pure Math.\ \textbf{33}, part 2,
Amer.\ Math.\ Soc., 1979, 313--346.

\bibitem{gamma:Gauss}
C.~F.~Gauss, \emph{Disquisitiones generales circa seriem infinitam}, Comment.\ Soc.\
Regiae Sci.\ Gottingensis Rec.\ \textbf{2} (1813), 3--46; Werke \textbf{3}, 265--327.
(The multiplication theorem.)

\bibitem{gamma:FT}
V.~Reshetnikov, \emph{FunctionTransformations} (Wolfram Language package), 2016,
\url{https://github.com/VladimirReshetnikov/FunctionTransformations}. (The $437$-rule
polylog and $670$-rule $\Gamma$ reduction tables anchoring the corpus.)

\bibitem{cm:ChowlaSelberg}
A.~Selberg and S.~Chowla, \emph{On Epstein's zeta-function}, J.~reine angew.\ Math.\
\textbf{227} (1967), 86--110.

\bibitem{cm:Damerell}
R.~M.~Damerell, \emph{$L$-functions of elliptic curves with complex multiplication, I},
Acta Arith.\ \textbf{17} (1970), 287--301.

\bibitem{cm:Hurwitz}
A.~Hurwitz, \emph{\"Uber die Entwicklungskoeffizienten der lemniskatischen Funktionen},
Math.\ Ann.\ \textbf{51} (1899), 196--226.

\bibitem{cm:Cox}
D.~A.~Cox, \emph{Primes of the Form $x^2+ny^2$: Fermat, Class Field Theory, and Complex
Multiplication}, Wiley, 1989. (Genus theory; class invariants; the $3\mid D$ obstruction
for $\gamma_2$.)

\bibitem{cm:MathWorldDigamma}
E.~W.~Weisstein, \emph{Digamma Function}, MathWorld --- A Wolfram Web Resource.
(Source of the FoxTrot identity.)

\bibitem{stieltjes:Adamchik98}
V.~S.~Adamchik, \emph{Polygamma functions of negative order}, J.~Comput.\
Appl.\ Math.\ \textbf{100} (1998), 191--199.

\bibitem{stieltjes:Apostol85}
T.~M.~Apostol, \emph{Formulas for higher derivatives of the Riemann zeta
function}, Math.\ Comp.\ \textbf{44} (1985), 223--232.

\bibitem{stieltjes:ApostolBook}
T.~M.~Apostol, \emph{Introduction to Analytic Number Theory},
Springer, 1976.

\bibitem{stieltjes:Berndt}
B.~C.~Berndt, \emph{Ramanujan's Notebooks, Part~I}, Springer, 1985.

\bibitem{stieltjes:BlagouchineRJ}
I.~V.~Blagouchine, \emph{Rediscovery of Malmsten's integrals, their
evaluation by contour integration methods and some related results},
Ramanujan J.\ \textbf{35} (2014), 21--110; Erratum--Addendum:
\textbf{42} (2017), 777--781.

\bibitem{stieltjes:CKK}
K.~Chakraborty, S.~Kanemitsu and T.~Kuzumaki, \emph{Finite expressions for
higher derivatives of the Dirichlet $L$-function and the Deninger
$R$-function}, Hardy--Ramanujan J.\ \textbf{32} (2009), 38--53.

\bibitem{stieltjes:Choudhury}
B.~K.~Choudhury, \emph{The Riemann zeta-function and its derivatives},
Proc.\ R.\ Soc.\ Lond.\ A \textbf{450} (1995), 477--499.

\bibitem{stieltjes:Coffey}
M.~W.~Coffey, \emph{Series representations for the Stieltjes constants},
Rocky Mountain J.\ Math.\ \textbf{44} (2014), 443--477
(arXiv:0905.1111).

\bibitem{stieltjes:Connon}
D.~F.~Connon, \emph{On some new formulae involving the Stieltjes constants},
arXiv:1902.00510 (2019).

\bibitem{stieltjes:EspinosaMoll04}
O.~Espinosa and V.~H.~Moll, \emph{A generalized polygamma function},
Integral Transforms Spec.\ Funct.\ \textbf{15} (2004), 101--115.

\bibitem{stieltjes:EspinosaMoll1}
O.~Espinosa and V.~H.~Moll, \emph{On some integrals involving the Hurwitz
zeta function: Part~1}, Ramanujan J.\ \textbf{6} (2002), 159--188.

\bibitem{stieltjes:EspinosaMoll2}
O.~Espinosa and V.~H.~Moll, \emph{On some integrals involving the Hurwitz
zeta function: Part~2}, Ramanujan J.\ \textbf{6} (2002), 449--468.

\bibitem{stieltjes:Kouba}
O.~Kouba, \emph{Generalized Stieltjes constants and integrals involving the
log--log function: Kummer's theorem in action}, J.~Classical Anal.\
\textbf{9} (2016), no.~2, 79--88 (arXiv:1702.05439).

\bibitem{stieltjes:MillerAdamchik}
J.~Miller and V.~S.~Adamchik, \emph{Derivatives of the Hurwitz zeta
function for rational arguments}, J.~Comput.\ Appl.\ Math.\ \textbf{100}
(1998), 201--206.

\bibitem{tower:BaileyBorwein2016}
D.~H.~Bailey and J.~M.~Borwein,
  \emph{Computation and structure of character polylogarithms with
  applications to character Mordell--Tornheim--Witten sums},
  Math.\ Comp.\ \textbf{85} (2016), 295--324.

\bibitem{Panzer2017}
E.~Panzer, \emph{The parity theorem for multiple polylogarithms},
J. Number Theory \textbf{172} (2017), 93--113.
\url{https://arxiv.org/abs/1512.04482}.

\bibitem{RivoalRoques}
T.~Rivoal and J.~Roques, \emph{Hadamard products of algebraic functions},
Section 3 (gamma relations and Rohrlich's conjecture).
\url{https://rivoal.perso.math.cnrs.fr/articles/algcoef.pdf}.

\bibitem{RadchenkoZagier}
D.~Radchenko and D.~Zagier, \emph{Arithmetic properties of the Herglotz function}.
\url{https://arxiv.org/abs/2012.15805}.
Author version: \url{https://people.mpim-bonn.mpg.de/zagier/files/preprints/HerglotzFunction.pdf}.

\bibitem{golden:AbouzahraLewin1985}
M.~Abouzahra and L.~Lewin, \emph{The polylogarithm in algebraic number fields},
Journal of Number Theory \textbf{21} (1985), no.~2, 214--244.
\href{https://doi.org/10.1016/0022-314X(85)90052-6}{doi:10.1016/0022-314X(85)90052-6}.

\bibitem{golden:AbouzahraLewinXiao1987}
M.~Abouzahra, L.~Lewin and H.~Xiao, \emph{Polylogarithms in the field of
omega (a root of a given cubic): Functional equations and ladders},
Aequationes Mathematicae \textbf{33} (1987), 23--45.
\href{https://doi.org/10.1007/BF01836149}{doi:10.1007/BF01836149}.

\bibitem{golden:BaileyBroadhurst1999}
D.~H.~Bailey and D.~J.~Broadhurst, \emph{A seventeenth-order polylogarithm
ladder}, preprint, 1999.
\url{https://arxiv.org/abs/math/9906134}.

\bibitem{golden:Gangl2013}
H.~Gangl, \emph{Functional equations and ladders for polylogarithms},
Communications in Number Theory and Physics \textbf{7} (2013), no.~3, 397--410.
\href{https://doi.org/10.4310/CNTP.2013.v7.n3.a1}{doi:10.4310/CNTP.2013.v7.n3.a1}.

\bibitem{CohenLewinZagier}
H.~Cohen, L.~Lewin and D.~Zagier, \emph{A Sixteenth-order Polylogarithm Ladder},
Experimental Mathematics \textbf{1} (1992), no. 1.
\url{https://people.mpim-bonn.mpg.de/zagier/files/exp-math-1/fulltext.pdf}.
(The higher-order ladders in this paper are explicitly conjectural numerical identities.)

\bibitem{HakimogluBrown}
C.~Hakimoglu-Brown, \emph{Analytic Dilogarithm Identities} (2025).
\url{https://arxiv.org/abs/2506.10206}.

\bibitem{RemiddiVermaseren}
E.~Remiddi and J.~A.~M.~Vermaseren, \emph{Harmonic polylogarithms},
Int. J. Mod. Phys. A \textbf{15} (2000), 725--754.
\url{https://arxiv.org/abs/hep-ph/9905237}.

\bibitem{Brown2012}
F.~Brown, \emph{Mixed Tate motives over $\mathbb Z$},
Annals of Mathematics \textbf{175} (2012), 949--976.
\url{https://arxiv.org/abs/1102.1312}.

\bibitem{zeros:DLMFGamma}
NIST Digital Library of Mathematical Functions, \S5.8, \emph{Infinite Products}; reciprocal-gamma Weierstrass product. \url{https://dlmf.nist.gov/5.8}. Accessed 7 October 2026.

\bibitem{hstruct:ZagierDilog}
D.~Zagier,
\emph{The dilogarithm function},
in P.~Cartier, B.~Julia, P.~Moussa, and P.~Vanhove (eds.),
\emph{Frontiers in Number Theory, Physics, and Geometry II},
Springer, 2007, pp.~3--65.
\url{https://doi.org/10.1007/978-3-540-30308-4_1}.

\bibitem{hstruct:Milnor}
J.~Milnor,
\emph{Introduction to Algebraic K-Theory},
Annals of Mathematics Studies \textbf{72},
Princeton University Press, 1971.
\url{https://doi.org/10.1515/9781400881796}.

\bibitem{hstruct:Washington}
L.~C.~Washington,
\emph{Introduction to Cyclotomic Fields}, second edition,
Graduate Texts in Mathematics \textbf{83}, Springer, 1997.
\url{https://doi.org/10.1007/978-1-4612-1934-7}.

\bibitem{hstruct:Sharifi}
R.~Sharifi,
\emph{Iwasawa theory}, online lecture notes, Chapter~4
(cyclotomic units and Dirichlet $L$-values).
\url{https://www.math.ucla.edu/~sharifi/notes/iwasawa-ch04.html}.
Accessed 8 October 2026.

\bibitem{hstruct:DLMF}
NIST,
\emph{Digital Library of Mathematical Functions},
\S\S2.11, 5.11, and 25.15; in particular the shifted Stirling
expansion and Dirichlet nonvanishing.
\url{https://dlmf.nist.gov/}.
Accessed 8 October 2026.

\bibitem{hstruct:Nemes}
G.~Nemes,
\emph{Error bounds for the asymptotic expansion of the Hurwitz
zeta function},
\emph{Proc. R. Soc. A} \textbf{473} (2017), 20170363.
\url{https://doi.org/10.1098/rspa.2017.0363};
\url{https://arxiv.org/abs/1702.05316}.

\bibitem{hstruct:Simon}
B.~Simon,
\emph{The classical moment problem as a self-adjoint finite
difference operator},
\emph{Adv. Math.} \textbf{137} (1998), 82--203.
\url{https://doi.org/10.1006/aima.1998.1728};
\url{https://arxiv.org/abs/math-ph/9906008}.

\bibitem{dist:Kubert}
D. S. Kubert, \emph{The universal ordinary distribution},
Bulletin de la Soci\'et\'e Math\'ematique de France
\textbf{107} (1979), 179--202.
\url{https://doi.org/10.24033/bsmf.1891}.

\bibitem{dist:Ouyang}
Y. Ouyang, \emph{On the universal norm distribution},
Journal of the Ramanujan Mathematical Society
\textbf{17} (2002), no.~4, 287--311.
\url{https://arxiv.org/abs/math/0210297}.

\bibitem{loggamma:EspinosaMoll2002b}
O. Espinosa and V. H. Moll, “On Some Integrals Involving the
Hurwitz Zeta Function: Part 2,” \emph{Ramanujan Journal}
\textbf{6} (2002), 449--468.
\url{https://www.math.tulane.edu/~vhm/papers_html/hurwitz2.pdf}.

\bibitem{loggamma:MilneANT}
J. S. Milne, \emph{Algebraic Number Theory}, version 3.08,
19 July 2020, especially Theorem 4.3 and
Proposition 7.55 with Remark 7.56.
\url{https://www.jmilne.org/math/CourseNotes/ANT.pdf}.

\bibitem{spectral:KowalskiNikeghbali}
E. Kowalski and A. Nikeghbali,
\emph{Mod-Poisson convergence in probability and number theory},
author manuscript.
\url{https://people.math.ethz.ch/~kowalski/mod-poisson.pdf}.
The permutation-cycle example explains the classical reciprocal-gamma
factor underlying the coefficient mechanism.

\bibitem{modular:RomikScherer}
D. Romik and R. Scherer,
\emph{Alternative summation orders for the Eisenstein series \(G_2\)
and Weierstrass \(\wp\)-function},
arXiv:1811.01523, version 2 (2020).
\url{https://arxiv.org/abs/1811.01523}.

\bibitem{modular:Ramanujan}
S. Ramanujan, \emph{On certain arithmetical functions},
Transactions of the Cambridge Philosophical Society
\textbf{22} (1916), no.~9, 159--184.
Transcription of the original paper:
\url{https://ramanujan.sirinudi.org/Volumes/published/ram18.html}.
\bibitem{XuWang2023}
C. Xu and W. Wang, \emph{Dirichlet type extensions of Euler sums},
Comptes Rendus Math\'ematique \textbf{361} (2023), 979--1010.
\url{https://doi.org/10.5802/crmath.453}.
The specialization used here is equation (20), section 3.4,
printed page 994.
\bibitem{gaussian:Radford}
D.~E.~Radford, A natural ring basis for the shuffle algebra and an
application to group schemes, \emph{Journal of Algebra} \textbf{58}
(1979), no.~2, 432--454.
\href{https://doi.org/10.1016/0021-8693(79)90171-6}{doi:10.1016/0021-8693(79)90171-6}.

\bibitem{gaussian:Umezawa}
R.~Umezawa, \emph{An explicit parity theorem for multiple polylogarithms},
arXiv:2508.02040, 2025.
\url{https://arxiv.org/abs/2508.02040}.

\bibitem{gaussian:Zagier}
D.~Zagier, Special values and functional equations of polylogarithms,
in L.~Lewin (ed.), \emph{Structural Properties of Polylogarithms},
Mathematical Surveys and Monographs \textbf{37}, American Mathematical
Society, Providence, RI, 1991, 377--400.
Author copy:
\url{https://people.mpim-bonn.mpg.de/zagier/files/tex/LewinPolylogarithms/fulltext.pdf}.

\bibitem{gaussian:NakamuraShiraishi}
H.~Nakamura and D.~Shiraishi, Landen's trilogarithm functional equation
and \(\ell\)-adic Galois multiple polylogarithms,
in M.~Morishita, H.~Nakamura, and J.~Ueki (eds.),
\emph{Low Dimensional Topology and Number Theory},
Springer Proceedings in Mathematics \& Statistics \textbf{456},
Springer, Singapore, 2025, 237--262.
\href{https://doi.org/10.1007/978-981-97-3778-9_8}{doi:10.1007/978-981-97-3778-9\_8}.
Author preprint: \url{https://arxiv.org/abs/2210.17182}.

\bibitem{gaussian:ACEMKM}
T.~Amdeberhan, M.~W.~Coffey, O.~Espinosa, C.~Koutschan,
D.~V.~Manna, and V.~H.~Moll, Integrals of powers of loggamma,
\emph{Proceedings of the American Mathematical Society}
\textbf{139} (2011), no.~2, 535--545.
\href{https://doi.org/10.1090/S0002-9939-2010-10589-0}{doi:10.1090/S0002-9939-2010-10589-0}.
Author preprint:
\url{https://www.math.tulane.edu/~vhm/papers_html/lg-subm.pdf}.

\bibitem{gaussian:BBB}
D.~H.~Bailey, D.~Borwein, and J.~M.~Borwein,
On Eulerian log-gamma integrals and Tornheim--Witten zeta functions,
\emph{The Ramanujan Journal} \textbf{36} (2015), nos.~1--2, 43--68.
\href{https://doi.org/10.1007/s11139-012-9427-1}{doi:10.1007/s11139-012-9427-1}.
Author preprint:
\url{https://www.davidhbailey.com/dhbpapers/log-gamma.pdf}.
The version-specific auxiliary sign check concerns the author copy dated
27 June 2012: \url{https://carmamaths.org/resources/jon/log-gamma.pdf},
Remark 3, equations (63)--(65). The separately linked July copy has different
equation numbering; the sign observation is not a journal-version audit.

\bibitem{gaussian:BorweinStraub2015}
J. M. Borwein and A. Straub, \emph{Relations for Nielsen polylogarithms},
Journal of Approximation Theory \textbf{193} (2015), 74--88.
\href{https://doi.org/10.1016/j.jat.2013.07.003}{doi:10.1016/j.jat.2013.07.003}.
\url{https://arminstraub.com/downloads/pub/nielsenrelations.pdf}.


\bibitem{gaussian:BBBL2001}
J. M. Borwein, D. M. Bradley, D. J. Broadhurst, and P. Lison\v ek,
\emph{Special values of multiple polylogarithms},
Transactions of the American Mathematical Society \textbf{353} (2001),
907--941.
\href{https://doi.org/10.1090/S0002-9947-00-02616-7}{doi:10.1090/S0002-9947-00-02616-7}.
\href{https://arxiv.org/abs/math/9910045}{arXiv:math/9910045}.


\bibitem{gaussian:VollingaWeinzierl2005}
J. Vollinga and S. Weinzierl, \emph{Numerical evaluation of multiple
polylogarithms}, Computer Physics Communications \textbf{167} (2005),
177--194.
\href{https://doi.org/10.1016/j.cpc.2004.12.009}{doi:10.1016/j.cpc.2004.12.009}.
\href{https://arxiv.org/abs/hep-ph/0410259}{arXiv:hep-ph/0410259}.

\bibitem{gaussian:Zhao}
J.~Zhao, \emph{Standard relations of multiple polylogarithm values at roots
of unity}, arXiv:0707.1459. \url{https://arxiv.org/abs/0707.1459}.
\bibitem{s4proof:hintze2021s4}
W.~Hintze, \emph{Closed form of the sum
$s_4=\sum_{n=1}^{\infty}(-1)^nH_n/(2n+1)^4$},
Mathematics Stack Exchange question~4160026 and answer by pisco,
3 June 2021.
\url{https://math.stackexchange.com/questions/4160026/}.

\bibitem{s4proof:mhzhao2020}
M.~H.~Zhao, \emph{MZV evaluation of hypergeometric series},
arXiv:2007.02508, version~10, 27 September 2020.
\url{https://arxiv.org/abs/2007.02508}.

\bibitem{s4proof:jqzhao2008}
J.~Zhao, \emph{Multiple polylogarithm values at roots of unity},
C.~R.~Math.~Acad.~Sci.~Paris \textbf{346} (2008), 1029--1032;
expanded version arXiv:0810.1064.
\url{https://arxiv.org/abs/0810.1064}.

\bibitem{s4proof:racinet2002}
G.~Racinet, \emph{Doubles m\'elanges des polylogarithmes multiples aux
racines de l'unit\'e}, Publications Math\'ematiques de l'IH\'ES
\textbf{95} (2002), 185--231.
\url{https://arxiv.org/abs/math/0202142}.

\bibitem{rigid:AartsFokkinkKruijtzer2001}
J.~Aarts, R.~Fokkink, and G.~Kruijtzer, Morphic numbers,
\emph{Nieuw Archief voor Wiskunde}, series~5, \textbf{2}, no.~1
(2001), 56--58.
\url{https://www.nieuwarchief.nl/serie5/pdf/naw5-2001-02-1-056.pdf}.

\bibitem{rigid:Ljunggren1960}
W.~Ljunggren, On the irreducibility of certain trinomials and
quadrinomials, \emph{Mathematica Scandinavica} \textbf{8} (1960),
65--70. \url{https://doi.org/10.7146/math.scand.a-10593}.

\bibitem{rigid:Tverberg1960}
H.~Tverberg, On the irreducibility of the trinomials
$x^n\pm x^m\pm1$, \emph{Mathematica Scandinavica} \textbf{8}
(1960), 121--126.
\url{https://doi.org/10.7146/math.scand.a-10599}.
Original scanned text:
\url{https://journals.msp.org/mscand/article/download/2744/2743/2775}.

\bibitem{rigid:BoydMartinThom2015}
D.~W.~Boyd, G.~Martin, and M.~Thom,
Squarefree values of trinomial discriminants,
\emph{LMS Journal of Computation and Mathematics} \textbf{18}
(2015), 148--169, especially Lemma~4.1(a),(b).
\url{https://doi.org/10.1112/S1461157014000436}.

\bibitem{saddle:DLMFLaplace}
NIST Digital Library of Mathematical Functions,
\S2.3(iii), Laplace's method.
\url{https://dlmf.nist.gov/2.3#iii}.

\bibitem{sharp:FS}
P.~Flajolet and R.~Sedgewick,
\emph{Analytic Combinatorics}, Cambridge University Press, 2009.
The smooth implicit-function schema and singularity transfer are the
classical background for the late-coefficient analysis.

\bibitem{jets:dlmf-dirichlet}
NIST Digital Library of Mathematical Functions, \S25.15,
\emph{Dirichlet $L$-functions}, including (25.15.4), (25.15.6),
and the nonvanishing statement (25.15.9). Accessed 9 October 2026.
\url{https://dlmf.nist.gov/25.15}.

\bibitem{cycloquot:Zhao}
J.~Zhao,
\emph{Multiple polylogarithm values at roots of unity},
C. R. Acad. Sci. Paris, Ser. I \textbf{346} (2008), 1029--1032.
\href{https://doi.org/10.1016/j.crma.2008.09.011}
{doi:10.1016/j.crma.2008.09.011}.
\url{https://www.numdam.org/item/10.1016/j.crma.2008.09.011.pdf}.
Its standard-relation framework includes more relations and depths
than the restricted presentation computed in this article.

\bibitem{hergjet:DLMF}
National Institute of Standards and Technology,
\emph{NIST Digital Library of Mathematical Functions},
Section~13.4, equation~13.4.4, and Section~13.8(iii),
equations~13.8.11--16.
\url{https://dlmf.nist.gov/13.4.E4};
\url{https://dlmf.nist.gov/13.8.iii}.
Consulted 9--10 October 2026.

\bibitem{hergjet:Temme}
N.~M.~Temme,
\emph{Uniform asymptotic expansions of a class of integrals in terms of
modified Bessel functions, with application to confluent hypergeometric
functions},
SIAM J. Math. Anal. \textbf{21} (1990), no.~1, 241--261.
\href{https://doi.org/10.1137/0521013}{doi:10.1137/0521013}.
% Insert these items inside the article's thebibliography environment.
\bibitem{cmnorm:Barquero}
A.~Barquero-Sanchez, L.~Cadwallader, O.~Cannon, T.~Genao, and R.~Masri,
\emph{Faltings heights of CM elliptic curves and special Gamma values},
Research in Number Theory \textbf{3} (2017), article 13.
\href{https://doi.org/10.1007/s40993-017-0077-7}{doi:10.1007/s40993-017-0077-7}.
Sections 3--5 give the analytic normalization of the Kronecker and
Chowla--Selberg formulas, including nonmaximal quadratic orders.

\bibitem{cmnorm:Cohen2026}
H.~Cohen, \emph{On the Argument of the Lerch, Chowla--Selberg Formula
and CM Values of $\eta(\tau)$}, arXiv:2606.02904v1 (2026).
\url{https://arxiv.org/html/2606.02904v1}.
Theorem 1.3 states the period-product formula with conductor correction;
Theorem 2.1 addresses the choice of representatives and eta phases.

\bibitem{cmnorm:Sutherland}
A.~V.~Sutherland, \emph{Computing Hilbert class polynomials with the
Chinese Remainder Theorem}, Mathematics of Computation \textbf{80}
(2011), 501--538.
\href{https://math.mit.edu/~drew/classpoly.html}{Author's implementation and references}.

\bibitem{cmnorm:YuiZagier}
N.~Yui and D.~Zagier, \emph{On the singular values of Weber modular
functions}, Mathematics of Computation \textbf{66} (1997), 1645--1662.
\href{https://people.mpim-bonn.mpg.de/zagier/files/doi/10.1090/S0025-5718-97-00854-5/fulltext.pdf}{Author-hosted full text}.

\bibitem{cmnorm:Stein}
W.~Stein, \emph{Modular Forms: A Computational Approach},
Graduate Studies in Mathematics, vol.~79, American Mathematical Society, 2007.
\url{https://wstein.org/books/modform/modform/level_one.html}.
\bibitem{cmgenus:Masri}
R.~Masri, \emph{Kronecker's solution of Pell's equation for CM fields},
Annales de l'Institut Fourier \textbf{63} (2013), 2287--2306.
\url{https://aif.centre-mersenne.org/item/10.5802/aif.2830.pdf}.
\bibitem{cmnonzero:Kohnen}
W.~Kohnen, \emph{Transcendence of zeros of Eisenstein series and other
modular functions}, Commentarii Mathematici Universitatis Sancti Pauli
\textbf{52} (2003), no.~1, 55--57.
\url{https://doi.org/10.14992/00008713}.
\bibitem{cmnonzero:ChiriacJorza}
L.~Chiriac and A.~Jorza, \emph{The variation of zeros of the Miller basis},
arXiv:2605.09731v1 (2026), introduction and Section~7.
\url{https://arxiv.org/abs/2605.09731}.


\bibitem{intdist:KubertSign}
D.~S. Kubert, \emph{The $\Z/2\Z$ cohomology of the universal ordinary
 distribution}, Bulletin de la Soci\'et\'e Math\'ematique de France
\textbf{107} (1979), 203--224.
\href{https://doi.org/10.24033/bsmf.1892}{doi:10.24033/bsmf.1892}.

\bibitem{intdist:DLMFHurwitz}
NIST Digital Library of Mathematical Functions,
\emph{Section 25.11: Hurwitz Zeta Function}, especially the defining series,
meromorphic continuation and multiplication formulas.
\url{https://dlmf.nist.gov/25.11}.


\bibitem{rw:AngularReport}
ProveIt research contribution, \emph{Gaussian Polylogarithms Beyond Parity:
Exact relation modules, an octahedral proof of the mixed harmonic sum,
and global angular-zero geometry}, October 2026.
Preserved source: \src{reports/angular-zero-continuation/article.tex}.

\bibitem{rw:DLMFDigamma}
NIST Digital Library of Mathematical Functions, \S5.9,
\emph{Integral Representations}, especially formula 5.9.13.
\url{https://dlmf.nist.gov/5.9}.

\bibitem{rw:DLMFPolylog}
NIST Digital Library of Mathematical Functions, \S25.12,
\emph{Polylogarithms}.
\url{https://dlmf.nist.gov/25.12}.

\bibitem{rw:DLMFHyperbolic}
NIST Digital Library of Mathematical Functions, \S4.36,
\emph{Infinite Products and Partial Fractions}, especially formula 4.36.3.
\url{https://dlmf.nist.gov/4.36}.

\bibitem{rw:LiuPego}
J.-G. Liu and R. L. Pego, \emph{On generating functions of Hausdorff
moment sequences}, Transactions of the American Mathematical Society
\textbf{368} (2016), 8499--8518. The author's arXiv version 4 includes
 a corrigendum. \url{https://arxiv.org/abs/1401.8052v4}.

\bibitem{envref:report}
ProveIt Contributors, \emph{Sharp Gaussian Bounds, Cyclic Distribution
Torsion, and Extremal Lerch Zeros}, research continuation, 10 October 2026.
Prepared against repository revision \src{a0a90ef318}.
Original source and certificates are preserved under
\src{reports/formal-reductions-continuation/continuations/extremal-cyclic-descent/}.


\bibitem{rattetra:2015}
V. Reshetnikov, \emph{A conjectured identity for tetralogarithms $\Li_4$},
Math.StackExchange question 1373123 (2015).
\url{https://math.stackexchange.com/questions/1373123/}.
\bibitem{rattetra:2017}
V. Reshetnikov, \emph{Several conjectured identities for polylogarithms},
MathOverflow question 261408 (2017).
\url{https://mathoverflow.net/questions/261408/}.
\bibitem{rattetra:report}
ProveIt Contributors, \emph{Polylogarithm Certificates and Sharp Euler Bounds},
research continuation, 10 October 2026, prepared against \src{4c173c06cc}.
Preserved under
\src{reports/uniform-transition-continuation/continuations/tetralogarithm-euler-polynomial-zeros/}.
\bibitem{lone:report}
ProveIt Contributors, \emph{Global Radial Monotonicity for Leading-Index-One
Double Polylogarithms}, research continuation, 10 October 2026,
prepared against \src{1539f353af}.
Preserved under
\src{reports/angular-zero-continuation/continuations/leading-one-global-radius/}.
\bibitem{latebudget:sharp}
ProveIt Contributors, \emph{Sharp Universal Euler Bounds},
research continuation, 10 October 2026.
Preserved under
\src{reports/real-order-threshold/continuations/sharp-universal-euler/}.
\bibitem{latebudget:golden}
ProveIt Contributors, \emph{Polylogarithm Research: Euler Bounds, Golden
Ladders and Cayley Reductions}, research continuation, 10 October 2026.
Preserved under
\src{reports/fractional-cayley-scaling/continuations/golden-cayley-double-turning/}.


\bibitem{ado:report}
ProveIt Contributors, \emph{Signed-Kernel Oscillation at Every Depth:
Minimal Pick Compensation, Zero Geometry, and Certified Polylogarithmic
Identities}, research continuation, 10 October 2026, prepared against
\src{3447bc59b78d138a7f0d536b66ac6e5cc5d5e49f}.
Original source and certificates are preserved under
\src{reports/angular-zero-continuation/continuations/all-depth-signed-oscillation/}.

\bibitem{stconv:Sun}
Z.-H. Sun, \emph{On the properties of invariant functions},
arXiv:2209.14625v1 (2022), Remark 3.2, printed p.~15.
\url{https://arxiv.org/abs/2209.14625}.
\bibitem{stconv:algebra-report}
ProveIt Contributors, \emph{A Convolution Algebra of Stieltjes Functions},
research continuation, 10 October 2026, prepared against \src{0efdbaa694}.
Preserved under \src{reports/stieltjes-correlation-calculus/continuations/convolution-algebra/}.
\bibitem{stconv:shifted-report}
ProveIt Contributors, \emph{Shifted Hurwitz-Zeta Correlations:
Exact Stieltjes Closure, Polygamma Contact Terms, and Log-Gamma Antiderivatives},
research continuation, 10 October 2026, observed source \src{0e5ab2f60f}.
Preserved under \src{reports/stieltjes-correlation-calculus/continuations/shifted-hurwitz-jets/}.
\bibitem{stconv:closure-report}
ProveIt Contributors, \emph{Bilinear Closure of Generalized Stieltjes Correlations},
research continuation, 10 October 2026. Individual source observations and
scope are retained in the original provenance note.
Preserved under \src{reports/stieltjes-correlation-calculus/continuations/bilinear-correlation-closure/}.
\bibitem{stconv:calculus-report}
ProveIt Contributors, \emph{A Convolution Calculus for Stieltjes Functions},
research continuation, 10 October 2026, prepared against \src{570b0567f3}.
Preserved under \src{reports/stieltjes-correlation-calculus/continuations/convolution-calculus/}.

\bibitem{ghharm:Gauss}
NIST Digital Library of Mathematical Functions, Eq.~15.4.20,
Gauss summation of the hypergeometric function at one.
\url{https://dlmf.nist.gov/15.4.E20}.

\bibitem{ghharm:Choi}
J.~Choi, Summation formulas involving binomial coefficients, harmonic
numbers, and generalized harmonic numbers, \emph{Abstract and Applied
Analysis} \textbf{2014} (2014), Article~501906, printed p.~6,
Conjectures~4 and~6 and Eqs.~(40)--(46).
\url{https://doi.org/10.1155/2014/501906}.

\bibitem{ghharm:report}
ProveIt Contributors, \emph{Exact Stieltjes and Harmonic Identities},
research continuation, 10 October 2026, prepared against
\nolinkurl{fc4d3bf80534ad7c901d3b8c9e71baf2df0064ed}.
Preserved under \src{reports/stieltjes-harmonic-resonance/continuations/stieltjes-harmonic-identities/}.


\bibitem{tw:DLMFLerch}
NIST Digital Library of Mathematical Functions, \S25.14, Lerch's transcendent.
\url{https://dlmf.nist.gov/25.14}.
\bibitem{tw:DLMFDirichlet}
NIST Digital Library of Mathematical Functions, \S25.15, Dirichlet $L$-functions.
\url{https://dlmf.nist.gov/25.15}.
\bibitem{tw:HuKim}
S.~Hu and M.-S.~Kim, On the Stieltjes constants and gamma functions with
respect to alternating Hurwitz zeta functions, \emph{Journal of Mathematical
Analysis and Applications} \textbf{509}(1) (2022), 125930.
\url{https://arxiv.org/abs/2106.14674}.
\bibitem{tw:ZhuHuKim}
H.~Zhu, S.~Hu and M.-S.~Kim, On the properties of alternating invariant
functions, arXiv:2505.10985v3 (2025), Theorem~1.13.
\url{https://arxiv.org/abs/2505.10985}.
\bibitem{tw:report}
ProveIt Contributors, \emph{Twisted Stieltjes Convolutions and Harmonic
Laurent Identities}, research continuation, 10 October 2026, prepared against
\nolinkurl{e0d9463bdee9685dfb1dddb819059cc738540c57}.
Preserved under \src{reports/stieltjes-correlation-calculus/continuations/twisted-harmonic-laurent/}.

\bibitem{xid:report}
ProveIt external research continuation,
\emph{Exact Identities: Collision Calculus, Shifted Harmonic Jets and
Weighted Arctangent Families}, 10 October 2026.
Original package retained at
\code{reports/stieltjes-correlation-calculus/continuations/collision-and-weighted-identities/}.
The weighted real evaluations are proved in Chapter~3; the remaining
claims retain separate intake and audit status. Global priority is not established.

\bibitem{pcc:reports}
ProveIt research continuations,
\emph{Periodic Stieltjes Collisions: Exact Contact Terms, Finite-Part Fubini
Laws, and Gamma--Polylogarithm Primitives} and
\emph{Periodic Stieltjes Contact Calculus: All-Index Collision Corrections,
Harmonic Derivative Anomalies, and Exact Zeta-Jet Integrals}, 10 October 2026.
Independent routes retained in the correlation spine as
\code{continuations/periodic-collision-contacts/} and
\code{continuations/periodic-stieltjes-contact/}; their immediate coincident
source is preserved as \code{continuations/coincident-stieltjes/}.
The fixed-coordinate periodic comparison is reconciled in Chapter~9;
the derivative anomaly was already established there.

\bibitem{pcc:DLMFGamma}
NIST Digital Library of Mathematical Functions, \S5.7,
\emph{Series Expansions of the Gamma Function}, especially formula~5.7.3.
\url{https://dlmf.nist.gov/5.7.E3}. Inspected 10 October 2026.

\bibitem{mx:report}
ProveIt research continuation,
\emph{Mixed Spectral Identities}, 10 October 2026,
prepared against \code{dcd95baeb7c0e914b138bddef6829c5b1d52b726}.
Preserved at \code{reports/stieltjes-harmonic-resonance/continuations/mixed-spectral-identities/}.
The mixed Gamma identity is derived from the classical Rogers--Dougall sum
in Chapter~8. Its further resonance, harmonic and collision claims retain
separate audit scopes. No global priority is asserted.

\bibitem{mx:DLMFDougall}
NIST Digital Library of Mathematical Functions, \S16.4,
\emph{Rogers--Dougall Very Well-Poised Sum}, formula~16.4.9.
\url{https://dlmf.nist.gov/16.4.E9}. Inspected 10 October 2026.

\bibitem{gf:Bounded}
ProveIt research contribution,
\emph{A certified obstruction in the weight-seven reduction search},
in the extremal cyclic-descent continuation,
\texttt{article/s6\_obstruction.tex}.
Pinned repository commit
\texttt{5a790187c8e186e41e2b990b4941cb7a1a3c7b6b}.
\href{https://github.com/VladimirReshetnikov/ProveIt/blob/5a790187c8e186e41e2b990b4941cb7a1a3c7b6b/Analysis/Polylogarithms/docs/reports/formal-reductions-continuation/continuations/extremal-cyclic-descent/article/s6_obstruction.tex}{Source article}.

\bibitem{gf:Depth}
ProveIt research contribution,
\emph{A Cayley normal form that preserves series depth},
in the polylogarithm--Stieltjes critical-transition continuation,
\texttt{article/05\_cayley\_depth.tex}.
Pinned repository commit
\texttt{5a790187c8e186e41e2b990b4941cb7a1a3c7b6b}.
\href{https://github.com/VladimirReshetnikov/ProveIt/blob/5a790187c8e186e41e2b990b4941cb7a1a3c7b6b/Analysis/Polylogarithms/docs/reports/uniform-transition-continuation/continuations/polylog-stieltjes-critical-transition/article/05_cayley_depth.tex}{Source article}.

\bibitem{gf:Radford}
D.~E.~Radford,
A natural ring basis for the shuffle algebra and an application to group
schemes, \emph{Journal of Algebra} \textbf{58} (1979), no.~2, 432--454.
\href{https://doi.org/10.1016/0021-8693(79)90171-6}{doi:10.1016/0021-8693(79)90171-6}.

\bibitem{gf:Patras}
F.~Patras,
La d\'ecomposition en poids des alg\`ebres de Hopf,
\emph{Annales de l'Institut Fourier} \textbf{43} (1993), no.~4, 1067--1087.
\href{https://doi.org/10.5802/aif.1365}{doi:10.5802/aif.1365}.


\bibitem{gf:report}
ProveIt research continuation, \emph{Exact Resonant Identities Beyond the Diagonal}, 11 October 2026. Original package preserved under \path{reports/stieltjes-harmonic-resonance/continuations/exact-resonant-identities/}. The complete-Cayley formal obstruction is audited in Chapter 12; the numerical $S_6$ conjecture remains open.

\bibitem{ghp:report}
ProveIt research continuation, \emph{Harmonic Parity and Resolvent Identities}, 11 October 2026. Original package preserved under \path{reports/stieltjes-harmonic-resonance/continuations/harmonic-parity-resolvents/}. The centered trigamma-square moments and their polylogarithmic generator are audited in Chapter 8; the broader classification and resolvent claims retain separate scopes.


\bibitem{dgp:reports}
ProveIt research continuations, \emph{Higher Digamma Powers and Ordered Gamma Sums},
\emph{Canonical Digamma Polynomials and the Fifth Harmonic Bridge}, and
\emph{Mixed Resolvent Products, Higher Digamma Powers, and Stieltjes Collisions},
11 October 2026. Original packages are preserved as
\path{reports/stieltjes-harmonic-resonance/continuations/higher-digamma-powers/},
\path{reports/stieltjes-harmonic-resonance/continuations/canonical-digamma-fifth/},
and \path{reports/stieltjes-correlation-calculus/continuations/mixed-resolvents-harmonic-powers/}.
Their all-power and fifth-moment cores are reconciled in Chapter 8; other
Gamma, resolvent, collision and ordered branches retain separate audit scopes.

\bibitem{dgp:DLMF}
NIST Digital Library of Mathematical Functions, Sections 5.5 (digamma
recurrence and reflection), 5.7 (Taylor and reciprocal-Gamma series), and
5.11 (sectorial asymptotics). In particular formulas 5.5.2, 5.5.4, 5.7.4
and 5.11.2. \url{https://dlmf.nist.gov/5}. Inspected 11 October 2026 UTC.

\bibitem{dgp:IKZ}
K. Ihara, M. Kaneko and D. Zagier, Derivation and double shuffle relations
for multiple zeta values, \emph{Compositio Mathematica} \textbf{142} (2006),
307--338. \url{https://doi.org/10.1112/S0010437X0500182X}.

\bibitem{dgp:Young}
P. T. Young, Series of height one multiple zeta functions, \emph{Integers}
\textbf{24} (2024), A43, published 20 May 2024.
\url{https://math.colgate.edu/~integers/y43/y43.pdf}.
\url{https://doi.org/10.5281/zenodo.11221606}.
The coefficient-description check in Appendix A concerns the inspected
publisher PDF, printed page 12; its ordinary Stieltjes formula is retained.

\bibitem{dgp:CanV1}
M. Can, L. Karg\i n, M. Cenkci and A. Dil, \emph{On certain harmonic zeta
functions}, arXiv:2403.07123v1, 11 March 2024, printed page 1.
\url{https://arxiv.org/pdf/2403.07123v1}.
The initial-index correction in Appendix A is scoped to that version's
introductory Stieltjes-limit display. Inspected 11 October 2026 UTC.


\bibitem{dpm:reports}
ProveIt research continuations, \emph{Bernoulli Determinants, Cubic Tornheim
Jets, and Canonical Digamma Integrals}, \emph{Differential Polygamma and
Lerch Identities}, and \emph{Exact Polylogarithm Parity, Harmonic-Tail
Reconstruction, and Higher Polygamma Moments}, 11 October 2026.
Original packages preserved under
\path{reports/stieltjes-harmonic-resonance/continuations/bernoulli-tornheim-digamma/},
\path{reports/stieltjes-harmonic-resonance/continuations/differential-polygamma-lerch-jets/},
and \path{reports/stieltjes-harmonic-resonance/continuations/polygamma-tail-identities/}.
Their common differential-polynomial analytic core is reconciled in
Chapter~8; their remaining coordinate, rank and Lerch
claims retain separate audit scopes.

\bibitem{btc:report}
ProveIt research continuation, \emph{Beta Mixtures Across Integer Twists:
Finite spectral corrections, harmonic logarithmic moments, and
Stieltjes--polylogarithm primitives}, 10 October 2026.
Original article, PDF and verification evidence preserved under
\path{reports/stieltjes-correlation-calculus/continuations/beta-crossing/}.
The affine crossing, integer-total mixture and balanced-jet cores are
reconciled in Chapter~8; the later logarithmic Mellin and primitive
families retain their separate integration scope.


\bibitem{opf:reports}
ProveIt research continuations, \emph{Alternating Tails and Spectral Defects},
\emph{Quartic Spectral Reconstruction}, \emph{Bernoulli Tails and Spectral Faces},
\emph{Bilinear Tails and Harmonic Determinants}, and
\emph{Tail Kernels and Boundary Products}, 11 October 2026.
Original sources and exact evidence preserved on the harmonic-resonance
spine and under \path{reports/tail-boundaries/}.
The common centered coefficient model, retained determinant and effective
eventual proof are reconciled in Chapter~8. Complete orders ten and twelve
retain report attribution; the collective order-fourteen extension has
its full exact certificate under \path{verification/prefix-order-fourteen/}.

\bibitem{ung:reports}
ProveIt research continuations, \emph{Universal Gamma Normal Forms:
Rational Kernels, Harmonic Sums, and Stieltjes Primitives}, and
\emph{A Universal Beta Formula for Gamma Normal Forms}, 11 October 2026.
Original articles, PDFs and exact evidence preserved under
\path{reports/stieltjes-harmonic-resonance/continuations/universal-gamma-coefficients/}
and \path{reports/stieltjes-harmonic-resonance/continuations/universal-beta-normal-forms/}.
Their common all-index coefficient formula and ordinary finite projections
are reconciled in Chapter~8. The collective analytic-jet extension gives
an actual locally convergent infinite normal generator.

\bibitem{ngp:reports}
ProveIt research continuations, \emph{Higher Identities},
\emph{Cubic Gamma Normal Form}, and \emph{Universal Gamma Coefficients},
11--12 October 2026. Delivered articles and exact evidence preserved on
the harmonic-resonance continuation spine. The compensator construction
uses the classical Euler homotopy operator: D.~Poole and W.~Hereman,
\emph{The homotopy operator method for symbolic integration by parts and
inversion of divergences with applications},
\url{https://arxiv.org/abs/0908.0399}; W.~Hereman, B.~Deconinck and
L.~D.~Poole, \emph{Continuous and Discrete Homotopy Operators:
A Theoretical Approach made Concrete},
\url{https://arxiv.org/abs/nlin/0510064}.
The hierarchy specialization, analytic primitive domain and independent
normal-generator proof are reconciled in Chapter~8. The collective
continuation proves every primitive coefficient and full support/sign, with
convergence on the full input disk; fixed-gap analytic completion remains
a separate pending proof scope.

\bibitem{rha:reports}
ProveIt research continuation, \emph{A Finite Polylogarithm Alphabet for
Rational Hankel Sums}, 11 October 2026. Original article, compiler and
recorded evidence preserved under
\path{reports/stieltjes-harmonic-resonance/continuations/rational-hankel-polylogarithms/}.
The confluent weight budget, anchored two-letter compiler and quadratic
denominator evaluation are reconciled in Chapter~8. The classical tools
retain their attribution: M.~E.~Hoffman, \emph{Quasi-shuffle products},
J.~Algebraic Combin.~\textbf{11} (2000), 49--68,
\url{https://arxiv.org/abs/math/9907173}; E.~Panzer,
\emph{Algorithms for the symbolic integration of hyperlogarithms with
applications to Feynman integrals}, Comput.~Phys.~Commun.~\textbf{188}
(2015), 148--166, \url{https://arxiv.org/abs/1403.3385}.

\bibitem{rj:reports}
ProveIt research continuations, \emph{Resonant Polylogarithmic Jets and
Convergent Inverse Transseries} and \emph{Polylogarithm Order Jets and
Resonant Inversion}, 11--12 October 2026.
Original sources and evidence are preserved under
\path{reports/uniform-transition-continuation/continuations/},
in \path{normalized-resonant-jets/} and
\path{order-jets-resonant-inversion/}.
The fixed-sector theorem extends R1 of the preserved
\path{resonance-correlations-nonseparable-jets/} report.
Reciprocal-Gamma normalization is credited to both overlapping
continuations. Section~\ref{rj:section} develops the exact order-lowering
and arbitrary repeated-primitive identities from that normalization.
Inverse geometry and Borel reconstruction retain separate review scopes.

\bibitem{qfb:reports}
ProveIt research continuations, \emph{Resonant Polylogarithm Transseries}
and \emph{Boundary Completion for Quadratic Polylogarithm Feedback},
11--12 October 2026. Original articles, PDFs and evidence are preserved
under \path{Analysis/Transseries/docs/series-and-transseries/}, in
\path{Resonant_Polylogarithm_Transseries/} and
\path{Quadratic_Polylogarithm_Feedback_Boundary_Completion/}.
The first report establishes the interior fold and poses the boundary
question; the second proves its completion. Section~\ref{qfb:section}
reconciles that classification and develops finite Bessel-polynomial
coefficient and order-spectrum identities. Nonmarginal critical charts and
crossover refinements retain separate integration scopes.

\bibitem{qfb:bessel}
NIST Digital Library of Mathematical Functions, \S18.34,
\emph{Bessel Polynomials}, especially (18.34.1)--(18.34.5),
\url{https://dlmf.nist.gov/18.34}. The convention here is
$y_n(x)=y_n(x;2)$.

\bibitem{qfb:transfer}
P.~Flajolet and A.~Odlyzko, \emph{Singularity analysis of generating
functions}, SIAM J.~Discrete Math.~\textbf{3} (1990), 216--240.

\bibitem{lmi:lambert}
NIST Digital Library of Mathematical Functions, \S4.13,
\emph{Lambert $W$-function}, \url{https://dlmf.nist.gov/4.13}.

\bibitem{lmi:pakes}
A.~G.~Pakes, \emph{Affine relation between an infinitely divisible
distribution function and its L\'evy measure}, J.~Math. Anal. Appl.\
\textbf{499} (2021), article 125008,
\url{https://doi.org/10.1016/j.jmaa.2021.125008}.
Related lower-branch Laplace-transform literature; the normalization
and identities used in Section~\ref{lmi:section} are proved there.

\bibitem{mtt:report}
ProveIt research continuation, \emph{Mixed Hurwitz-Tail Moments,
Zeta-Anchor Reductions, and Stieltjes Transport}, 11 October 2026,
preserved article and verification evidence in
\path{reports/stieltjes-harmonic-resonance/continuations/mixed-tail-anchor-reductions/}.
Its third research question asks for rational polynomial combinations
for $(2,2,3)$ and $(2,3,4)$. Section~\ref{mtt:section} develops the first
case within a precisely specified polynomial certificate class.

\bibitem{llc:inverse}
ProveIt research continuation, \emph{Double Resonance and Nested Stokes
Geometry in Polylogarithmic Lambert Inversion}, 12 October 2026.
Preserved source and evidence in
\path{Analysis/Transseries/docs/series-and-transseries/Resonant_Lambert_Inverse_Transseries/}.
The finite inverse formula and the conjectured hyperbolic amplitudes
motivate Section~\ref{llc:section}.

\bibitem{llc:reversion}
ProveIt research continuation, \emph{Summation and Singularities under
Polylogarithmic Reversion}, 12 October 2026, in the preceding report's
\path{continuations/summation-singularities/} directory.
Its Lambert section proves the leading amplitudes and first correction;
its all-order question is resolved by Theorem~\ref{llc:all-orders}.
Original source evidence and fresh isolated verification are preserved.

\end{thebibliography}
